\paragraph*{Problems}

\begin{enumerate}[label=\thesection.\arabic*.,ref=\thesection.\arabic*,start=20,series=book-problems]

\item\label{ex:piecewise-smooth-distribution-derivative} 设 $f$ 为具有分片连续导数的分片连续函数. 描述 $f$ 的分布导数.

\item\label{ex:logarithm-distribution-derivative} 求函数 $\ln|x|$ 的分布导数.

\item\label{ex:riemann-integrable-traveling-wave} 设 $u(x,t)=f(x+t)$, 其中 $f$ 为 $\mathbb{R}$ 上任意局部 Riemann 可积函数. 证明在分布意义下 $u_{tt}=u_{xx}$.

\item\label{ex:inverse-square-distribution-laplacian} 在 $\mathbb{R}^3$ 中计算 $\Delta(1/r^2)$.

\item\label{ex:exponential-oscillation-tempered} 证明 $e^x\cos e^x\in\mathcal{S}^{\prime}(\mathbb{R})$.

\item\label{ex:polynomial-cosine-coefficients-distribution} 证明, 只要当 $n\to\infty$ 时 $|a_n|$ 至多以多项式速度增长, 则 $\sum_{n\in\mathbb{N}}a_n\cos nx$ 在分布意义下收敛.

\item\label{ex:prove-delta-primitive-conditions} 补全 Example~\ref{ex:delta-sequence-primitive} 的细节.

\item\label{ex:distribution-ode-system-substitution} 讨论如何将代换 \eqref{eq:distribution-ode-primitive-reduction} 推广到常微分方程组.

\item\label{ex:singular-ode-heaviside-general-solution} 证明 $xy'=0$ 的通解为 $c_1+c_2H(x)$. Hint: 先证明若 $\phi\in\mathcal{D}(\mathbb{R})$ 在原点为零, 则 $\phi(x)/x$ 为试验函数.

\item\label{ex:translation-invariant-distribution-constant} 设 $f\in\mathcal{D}^{\prime}(\mathbb{R})$ 对每个正的 $h$ 都满足 $f(x+h)=f(x)$. 证明 $f$ 为常数.

\item\label{ex:distribution-primitives-convergent-subsequence} 设 $f_n$ 为 $\mathcal{D}^{\prime}(\mathbb{R})$ 中的收敛序列, 并假设 $F_n'=f_n$. 再假设存在积分非零的试验函数 $\phi_0$, 使序列 $(F_n,\phi_0)$ 有界. 证明 $F_n$ 有一个收敛子列.

\item\label{ex:even-distribution-odd-primitive} 证明 $\mathbb{R}$ 上的偶分布具有奇原函数.

\item\label{ex:point-supported-distribution-delta-derivatives} 假设分布 $f$ 的支集为集合 $\{0\}$. 证明 $f$ 为 delta 函数的导数的线性组合. Hint: 取 $n$ 如 Lemma~\ref{lem:distribution-local-finite-order} 所给出, 并假设当 $|\alpha|\leq n$ 时 $D^\alpha\phi(0)$ 为零. 设 $e$ 为试验函数, 当 $|\bm{x}|\leq1$ 时等于 1, 当 $|\bm{x}|\geq2$ 时等于 0. 现在考虑序列 $\phi_k(\bm{x})=\phi(\bm{x})e(k\bm{x})$. 证明 $(f,\phi_k)\to0$, 从而 $(f,\phi)=0$.

\end{enumerate}
